\documentclass{article}
\usepackage{amsmath}
\usepackage{graphicx}
\usepackage{circuitikz}

% FWC IITH – LaTeX: GATE IN 2022 Question Paper (Digital Circuits)
% Author: Tirumala Sai Nithin | FWC22187 | IIT Hyderabad
% Toolchain: pdflatex via Termux Debian (Neovim on Android)
% Build: pdflatex gate_in_2022.tex → termux-open gate_in_2022.pdf

\newcommand{\figref}[1]{Fig.~\ref{#1}}

\begin{document}

\begin{center}
\textbf{GATE - IN 2022}\\
TIRUMALA SAI NITHIN\\
December 2023
\end{center}

\begin{enumerate}

\item The digital circuit shown in \figref{fig:GATEIN202236}
\begin{figure}[h]
\centering
\includegraphics[width=\columnwidth]{figs/GATEIN202236.png}
\caption{GATE IN 2022 Digital Circuit}
\label{fig:GATEIN202236}
\end{figure}
\begin{enumerate}
\item is a divide-by-$5$ counter
\item is a divide-by-$7$ counter
\item is a divide-by-$8$ counter
\item does not function as a counter due to disjoint cycles of states
\end{enumerate}
\hfill(GATE-IN-2022)

\end{enumerate}

\textbf{Solution:}\\
The circuit implements a 3-bit LFSR-style shift register with feedback:
\[
Q_0^{n+1} = \overline{Q_2^n}, \quad Q_1^{n+1} = Q_0^n, \quad Q_2^{n+1} = Q_1^n
\]
State sequence: $000 \to 001 \to 011 \to 111 \to 110 \to 100 \to 000$\\
Five unique states $\Rightarrow$ \textbf{Divide-by-5 counter (Option A)}

\end{document}
